\section{Introduction}
Heterogeneous treatment effects motivate analysis across the potential-outcome distribution, but endogenous treatment prevents conventional quantile regression from recovering structural effects. Instrumental variable quantile regression (IVQR) combines excluded instruments with structural-rank restrictions to address this problem \parencite{koenker_bassett_1978,chernozhukov_hansen_2005}. In applications with many controls, nuisance estimation introduces an additional difficulty. Regularization and Neyman-orthogonal moments can reduce its impact on inference, but neither creates identifying information when instruments explain little treatment variation.

The IVQR literature distinguishes moment validity from identification and informative inference \parencite{chernozhukov_hansen_wuthrich_2017}. Inverse quantile regression provides practical estimation under regular identification \parencite{chernozhukov_hansen_2006}; candidate-specific test inversion accommodates weak identification under its maintained regularity conditions \parencite{chernozhukov_hansen_2008}. Its accepted regions can nevertheless be wide or disconnected. The finite-sample conditional procedure of \textcite{chernozhukov_hansen_jansson_2009} is a separate inferential construction. Related developments in the broader IV literature include the direct structural tests of \textcite{anderson_rubin_1949}, weak-IV approximations \parencite{staiger_stock_1997}, weak-identification-robust GMM \parencite{stock_wright_2000}, and pivotal and conditional likelihood-ratio procedures \parencite{kleibergen_2002,moreira_2003}. \textcite{andrews_stock_sun_2019} review their implications. Conventional first-stage diagnostics describe relevance, but do not provide universal quantile-specific IVQR identification thresholds.

High-dimensional regularization addresses a different source of difficulty. \textcite{belloni_chernozhukov_2011} develop penalized quantile regression and pivotal penalty calibration; \textcite{belloni_chernozhukov_kato_2019} establish orthogonal-score inference in approximately sparse quantile models. Double/debiased machine learning (DML) uses Neyman orthogonality to reduce first-order sensitivity to sufficiently small nuisance errors \parencite{chernozhukov_et_al_2018}. Cross-fitting is a distinct device for separating nuisance estimation from score evaluation. \textcite{chen_huang_tien_2021} combine regularized nuisance estimation and an orthogonal IVQR score, comparing Oracle-GMM, Full-GMM, and DML-IVQR in a high-dimensional Monte Carlo design. Further approaches address endogenous binary treatment \parencite{wuthrich_2019} and orthogonal moments for nonparametric IVQR with high-dimensional covariates \parencite{jin_sun_2025}.

The high-dimensional simulation in \textcite{chen_huang_tien_2021} holds instrument relevance fixed. We ask how progressively weakening that relevance changes point accuracy, coverage, power, and confidence-set informativeness across quantiles. Our contribution is a controlled Monte Carlo extension of their design and preserved executable procedures. A relevance parameter scales the excluded signal, while compensating independent noise preserves the treatment distribution and its joint distribution with the outcome disturbance. The extension introduces neither a new estimator nor a new inference theory.

The experiment compares three complete implementations on common samples. Oracle-GMM knows the designated relevant controls; Full-GMM uses all controls without regularization; DML-IVQR regularizes nuisance estimation. As relevance falls, all three lose point accuracy and numerical localization, while power generally declines. DML-IVQR usually improves point accuracy relative to Full-GMM, with exceptions. Coverage can approach 95\% as tests lose discrimination, and several strong-relevance cells exhibit substantial undercoverage. The implementation differences, non-finite population density weight, and asymptotic reference distribution qualify interpretation throughout. These results separate the benefits of nuisance estimation from the information required to distinguish competing structural effects.
